\section{ML Pipeline Design}
\label{sec:pipeline-design}

All pipelines share candidate generation, splits, the scorer and its
hyperparameters, and the feature construction of
\S\ref{sec:problem-setup}. Only the encoder that produces
$\mathbf{e}_i$ differs.

\subsection{Encoders: a capacity ladder}
\label{sec:baseline-pipelines}
\label{sec:large-embedding-pipeline}
\label{sec:llm-augmented-pipelines}
\label{sec:hybrid-pipelines}
We compare six item representations in increasing order of capacity and
cost:
\begin{itemize}
  \item \textbf{TF-IDF}: a 2{,}000-term TF-IDF vector over the item
    text, which is classical and involves no model inference.
  \item \textbf{SBERT}: \texttt{all-MiniLM-L6-v2} \cite{sbert}, a
    22M-parameter sentence encoder. It is the cheapest dense encoder
    and the reference point for every cost comparison.
  \item \textbf{mxbai}: \texttt{mxbai-embed-large}, a 335M-parameter
    dedicated embedding encoder (BERT-style, not an LLM), served
    through Ollama.
  \item \textbf{phi4-mini (raw)}: the 3.8B-parameter decoder-only chat
    LLM \texttt{phi4-mini}, used as an embedder by mean-pooling its
    last-layer hidden states with no embedding-specific training
    (\texttt{llama-server --pooling mean}). This is what a practitioner
    gets by pointing an existing chat model at an embedding endpoint.
  \item \textbf{Qwen3-Embedding-4B (trained)}: a 4B-parameter
    decoder-only LLM \emph{contrastively trained for embedding}, served
    with last-token pooling as its model card specifies. Its size and
    quantization (Q4\_K\_M) match phi4-mini's, so comparing the two
    separates the effect of LLM capacity from the effect of using
    untrained hidden states as an embedding
    \cite{sgpt, llm2vec, e5-mistral}.
  \item \textbf{Hybrid}: TF-IDF concatenated with phi4-mini. The two
    vectors each contribute their own affinity features.
\end{itemize}
SBERT and mxbai accept at most 256 and 512 input tokens respectively,
while review-based item text runs longer. We therefore split long text
into word chunks that fit each encoder (150 and 300 words), embed each
chunk, and mean-pool the chunk vectors, so that every encoder sees the
whole item text. Both LLMs run with a 4{,}096-token context, which
covers every item without chunking. All models run locally on one Apple
M4 machine (16\,GB), and all GGUF models use the same Q4\_K\_M
quantization.

\subsection{Reference baselines}
Two baselines use no item text at all and give the absolute numbers a
reference point. \textbf{Popularity} ranks candidates by $\log(1+n_i)$.
\textbf{PureSVD} \cite{puresvd} factorizes the binarized training
user--item matrix (rank 50, fixed in advance rather than tuned on test
lists) and scores $(u,i)$ by the reconstructed entry. A content pipeline
that cannot beat popularity is not using its item representation for
personalization. A pipeline that cannot beat PureSVD gains nothing from
text over a standard collaborative model.

\subsection{Scorers}
We use two scorers so that conclusions are not tied to one model's
inductive biases. The first is LightGBM's \texttt{LGBMRanker}
\cite{lightgbm} (lambdarank objective, grouped by user, 200 trees, 64
leaves, learning rate 0.05). The second is a pointwise multi-layer
perceptron (scikit-learn \texttt{MLPRegressor}, hidden layers 64 and 32,
standardized inputs, early stopping) that regresses the graded
relevance label. Hyperparameters are fixed across all pipelines and
datasets.
